% Converted from sample_problem_set_source.pdf by an AI agent, in one
% pass from the PDF alone (no source .tex was available). The solution
% boxes are filled in with worked answers to show what a finished
% submission would look like.

\documentclass[11pt]{article}

\usepackage{amsmath, amssymb}
\usepackage[most]{tcolorbox}
\usepackage[margin=1in]{geometry}
\usepackage{parskip}

% Solution box: blue-tinted, bold "Solution:" lead-in, roughly matching
% the appearance of the boxes in the source PDF.
\newtcolorbox{solution}{
  colback=blue!4!white,
  colframe=blue!45!black,
  boxrule=0.5pt,
  arc=1.5pt,
  left=6pt, right=6pt, top=5pt, bottom=5pt,
  before upper={\textbf{Solution:}\ },
}

\title{Problem Set 0: Math and Statistics Review \\
       \large USC Marshall --- Financial Derivatives FBE 459}
\author{Scott Joslin}
\date{Spring 2026 \\ (Due: Thursday January 22, 11:59 PM)}

\begin{document}
\maketitle

\noindent\emph{If possible, please write your answer on a printed copy of
the problem set and submit a scan/photograph of your solutions. Show your
work cleanly and write out the formulas that you used to solve the
problems. Circle your final answers. You may collaborate freely with your
classmates. You may freely consult AI such as ChatGPT to assist you.
However, you should not ask it to directly solve the problem and you must
write your answer in your own words. Please give credit to any source of
assistance (reference, person, AI, etc.) with the assignment.}

\medskip

\noindent\emph{This problem set is graded on completeness. Review the
posted solutions to check your answers. These topics will be revisited
later in the course.}

\section*{Problem 1}

\textbf{Derivatives.} Suppose that $W(t)$ represents your wealth at time
$t$. Say at that your wealth has the derivative:
\begin{equation}
  W'(t) \equiv \frac{dW}{dt} = 0.12 \times W(t). \tag{$*$}
\end{equation}

\begin{enumerate}
  \item[(a)] If time is measured in years and wealth is measured in
  dollars, what are the units of $W'(t)$?
  \begin{solution}
  Dollars per year (\$/year). A derivative has units of the numerator
  divided by the denominator: $dW$ is in dollars, $dt$ is in years.
  \end{solution}

  \item[(b)] In english words, what does the derivative $dW/dt$ mean?
  \begin{solution}
  The instantaneous rate at which wealth is changing with respect to
  time --- i.e., how fast wealth is growing (or shrinking) at that exact
  moment, in dollars per year.
  \end{solution}

  \item[(c)] Recall that the derivative is the instantaneous rate of
  change. For a small time interval, $\Delta t > 0$, we can approximate
  the change in wealth by:
  \begin{equation*}
    \frac{\Delta W}{\Delta t} \approx W'(t_0) = 0.12 \times W(t_0).
  \end{equation*}
  Here, $\Delta W = W(t_1) - W(t_0)$ and $\Delta t = t_1 - t_0$. Based on
  this approximation, if you start with \$100 at $t = 0$ (so that $W(0) =
  \$100$), approximately what would $W(1/12)$ be?
  \begin{solution}
  Using the approximation with $W(0) = \$100$, $\Delta t = 1/12$:
  \begin{equation*}
    \Delta W \approx 0.12 \times W(0) \times \Delta t
                   = 0.12 \times 100 \times \tfrac{1}{12}
                   = \$1.
  \end{equation*}
  So $W(1/12) \approx 100 + 1 = \boxed{\$101}$.
  \end{solution}
\end{enumerate}

\noindent\textbf{Note:} you won't be asked to do things like: find the
derivative of $f(x) = 3x^2$; you will be expected to use concepts like
these.

\section*{Problem 2}

\textbf{Partial Derivatives.} A common approach to for stock valuation
that you probably learned in your first finance class is the Gordon
growth model. Under certain assumptions, this model states that:
\begin{equation*}
  P_0 = \frac{D_1}{r_{cc} - g}
\end{equation*}
where $P_0$ is the current price of the stock, $D_1$ is the expected
dividend at $t = 1$, $r_{cc}$ is the cost of capital, and $g$ is the
expected growth rate of dividends. Recall that the cost of capital is
the risk-adjusted discount rate:
\begin{equation*}
  r_{cc} = r_{\text{risk-free}} + (\text{premium}).
\end{equation*}
This formula expresses the price as a function of three variables:
$(D_1, r_{cc}, g)$.

\begin{enumerate}
  \item[(a)] In english words, what does the partial derivative
  $\dfrac{\partial P_0}{\partial r_{cc}}$ represent?
  \begin{solution}
  The sensitivity of the current stock price to a small change in the
  cost of capital, \emph{holding the expected dividend $D_1$ and growth
  rate $g$ fixed}. Concretely: how many dollars the price moves per unit
  change in $r_{cc}$.
  \end{solution}

  \item[(b)] Should the partial derivative $\dfrac{\partial P_0}{\partial r_{cc}}$
  be positive or negative?
  \begin{solution}
  Negative. A higher discount rate makes future dividends worth less
  today, so the price falls. (You can also verify algebraically: the
  denominator $r_{cc} - g$ grows with $r_{cc}$, shrinking $P_0$.)
  \end{solution}

  \item[(c)] Suppose that $D_1 = \$1$, $r_{cc} = 10\%$, and $g = 5\%$ so
  that $P = \$20$. Suppose that you are given that
  $\dfrac{\partial P}{\partial r_{cc}} = -\$400$. If $r_{cc}$ increases
  from 10\% to 10.1\% (an increase of $\Delta r_{cc} = 0.01$),
  approximately how much should the price change by if the dividend
  level and growth rate stay the same?
  \begin{solution}
  Using the linear approximation,
  \begin{equation*}
    \Delta P \approx \frac{\partial P}{\partial r_{cc}} \cdot \Delta r_{cc}
             = -\$400 \times 0.001 = -\$0.40.
  \end{equation*}
  (Here $\Delta r_{cc} = 0.001$ corresponds to the $10\% \to 10.1\%$
  move, i.e.\ a $0.1$ percentage-point increase written in decimal
  form.) So the new price is approximately $P \approx \$20 - \$0.40
  = \boxed{\$19.60}$.
  \end{solution}
\end{enumerate}

\noindent\textbf{Note:} For part (b), you could compute a formula using
rules from calculus. The question is asking you to think intuitively
without actually taking the derivative.

\section*{Problem 3}

\textbf{Expectation and Volatility.} Suppose that a certain stock will
either have a return of $+20\%$ or $-10\%$. There are no other
possibilities. There is a 60\% (40\%) probability that the return will be
$+20\%$ ($-10\%$).

What is the expected return of the stock?

\begin{solution}
\begin{equation*}
  E[R] = 0.6 \times 20\% + 0.4 \times (-10\%) = 12\% - 4\% = \boxed{8\%}.
\end{equation*}
\end{solution}

\section*{Problem 4}

\textbf{Normal distribution.}

\begin{enumerate}
  \item[(a)] \textbf{(10 points)} Suppose that the random variable $Z$
  follows the standard normal distribution. That is, $Z$ has a normal
  distribution with mean 0 and variance of 1. Find $P(Z < 1)$. You will
  need a standard normal table or program such as Excel.
  \begin{solution}
  From the standard normal CDF (or \texttt{=NORM.S.DIST(1, TRUE)} in
  Excel):
  \begin{equation*}
    P(Z < 1) = \Phi(1) \approx \boxed{0.8413}.
  \end{equation*}
  \end{solution}

  \item[(b)] \textbf{(10 points)} Suppose that the a certain stock return
  has a normal distribution with mean 10\% and volatility of 30\%. Find
  the probability that the stock return is less than 40\%. To find this,
  first compute the $Z$-score:
  \begin{equation*}
    Z = \frac{R - \mu_R}{\sigma_R}
  \end{equation*}
  \begin{solution}
  With $\mu_R = 10\%$, $\sigma_R = 30\%$, and $R = 40\%$:
  \begin{equation*}
    Z = \frac{40\% - 10\%}{30\%} = 1,
  \end{equation*}
  so
  \begin{equation*}
    P(R < 40\%) = P(Z < 1) = \Phi(1) \approx \boxed{0.8413}.
  \end{equation*}
  \end{solution}
\end{enumerate}

\end{document}
